\documentclass[12pt]{article}

\usepackage[a4paper,margin=1in]{geometry}
\usepackage[T1]{fontenc}
\usepackage[utf8]{inputenc}
\usepackage{lmodern}
\usepackage{booktabs}
\usepackage{longtable}
\usepackage{array}
\usepackage{enumitem}
\usepackage{hyperref}
\usepackage{setspace}
\usepackage{ragged2e}

\setstretch{1.08}
\setlength{\parskip}{0.5em}
\setlength{\parindent}{0pt}

\hypersetup{
	colorlinks=true,
	linkcolor=black,
	citecolor=black,
	urlcolor=blue,
	pdftitle={Verification Package Overview for the General Theory of Cognitive Structuring},
	pdfauthor={Kostiantyn Osmolovskyi}
}


\title{Verification Package Overview\\
	\small for the General Theory of Cognitive Structuring}

\author{
	Kostiantyn Osmolovskyi\\
	\small Independent Researcher, Ukraine\\
	\small ORCID: 0009-0006-3144-7237\\
	\small \texttt{constantinosmol@gmail.com}
}

\date{}

% ---------------------------------------------------------------------------
% Part of a series: General Theory of Cognitive Structuring (GTCS)
% Autor: Kostiantyn Osmolovskyi (ORCID: 0009-0006-3144-7237)
% License: Creative Commons Attribution 4.0 International (CC BY 4.0)
% Repository: https://github.com/k-osmolovskyi/k-osmolovskyi.github.io
% DOI: https://doi.org/10.5281/zenodo.19705773
% ---------------------------------------------------------------------------

\begin{document}
	\maketitle

\begin{center}
	\small
	Verification Report No. 2026-4\\
	Version 1.0
\end{center}

\vspace{0.5em}
	
	\section{Purpose of this document}
	
	This document provides a short overview of the public verification package accompanying the
	\emph{General Theory of Cognitive Structuring}. Its purpose is to make the structure of the
	package explicit, to clarify the role of each document, and to reduce entry friction for
	external readers.
	
	The package is intended to support external verification by providing separate but coordinated
	layers for:
	\begin{itemize}[leftmargin=2em]
		\item terminology,
		\item dependency logic,
		\item structural traceability,
		\item formal traceability,
		\item procedural entry,
		\item review practice,
		\item claims stabilization,
		\item and proof-status control.
	\end{itemize}
	
	This overview does not restate the theory itself. It serves as a front document for the
	verification infrastructure built around the series.
	
	\section{What the verification package is for}
	
	The verification package exists to make external reading of the theory easier, more accurate,
	and less vulnerable to accidental misinterpretation.
	
	More specifically, it is intended to:
	\begin{itemize}[leftmargin=2em]
		\item reduce category errors in the interpretation of core terms;
		\item clarify how the paper series unfolds without circular support;
		\item make the dependency structure of the series explicit;
		\item make the formal relation among notation, assumptions, definitions, and results traceable;
		\item provide procedural entry points for external readers;
		\item provide compact tools for review and checking;
		\item distinguish main claims from local commentary;
		\item and identify which results are already proof-ready and which remain pending later consolidation.
	\end{itemize}
	
	The package is therefore not a substitute for the primary papers. It is a structured support
	layer for verifying how the series is built and how its main components relate.
	
	\section{Documents included in the package}
	
	\begin{longtable}{>{\RaggedRight\arraybackslash}p{0.34\textwidth}
			>{\RaggedRight\arraybackslash}p{0.56\textwidth}}
		\toprule
		\textbf{Document} & \textbf{Role} \\
		\midrule
		\endfirsthead
		\toprule
		\textbf{Document} & \textbf{Role} \\
		\midrule
		\endhead
		
		\emph{Glossary of Core Terms} & Interpretive glossary for the main terms of the series. Stabilizes terminology and reduces accidental misreading. \\
		
		\emph{Acyclicity Statement and Dependency Criteria} & Public statement of the dependency criteria under which the series is non-circular. Distinguishes logical dependency from cross-reference, refinement, and synthesis. \\
		
		\emph{Dependency Tables and Node Registry} & Concrete node registry and paper-level dependency map showing what depends on what across the series. \\
		
		\emph{Technical Appendix} & Formal traceability layer for notation, assumptions, definitions, operators, constructions, and result-level dependencies. \\
		
		\emph{How to Verify the Theory} & Procedural roadmap for external readers. Explains reading order and what each companion document is for. \\
		
		\emph{External Verification Checklist} & Practical review instrument organized by verification task. Supports structured external checking. \\
		
		\emph{Minimal Claims Register} & Compact register of the theory's main claims, with claim type, dependencies, scope, source, and status. \\
		
		\emph{Collected Results and Proof Status Note} & Canonical author-side collected layer of the main results of the series together with proof-status classification. \\
		
		\emph{Collected Propositions and Theorems} & Normalized collection of proposition- and theorem-level results intended to support later proof consolidation. \\
		
		\emph{Proof Status Register} & Compact register distinguishing proof-ready results, sketch-level results, canonical synthetic results, and priority targets for future proof-compendium treatment. \\
		
		\bottomrule
	\end{longtable}
	
	\section{Recommended reading order}
	
	For a reader approaching the verification package for the first time, the following order is recommended:
	
	\begin{enumerate}[leftmargin=2em]
		\item \emph{Glossary of Core Terms}
		\item \emph{Acyclicity Statement and Dependency Criteria}
		\item \emph{Dependency Tables and Node Registry}
		\item \emph{Technical Appendix}
		\item \emph{How to Verify the Theory}
		\item \emph{External Verification Checklist}
		\item \emph{Minimal Claims Register}
		\item \emph{Collected Results and Proof Status Note}
		\item \emph{Collected Propositions and Theorems}
		\item \emph{Proof Status Register}
	\end{enumerate}
	
	This order moves from interpretive stabilization to structural verification, then to formal
	traceability, then to claim and result control, and finally to proof-status organization.
	
	\section{Alternative reading routes}
	
	Different readers may need different entry paths.
	
	\subsection*{For terminology and conceptual clarity}
	\begin{enumerate}[leftmargin=2em]
		\item \emph{Glossary of Core Terms}
		\item \emph{How to Verify the Theory}
		\item selected primary papers
	\end{enumerate}
	
	\subsection*{For dependency and non-circularity questions}
	\begin{enumerate}[leftmargin=2em]
		\item \emph{Acyclicity Statement and Dependency Criteria}
		\item \emph{Dependency Tables and Node Registry}
		\item selected primary papers
	\end{enumerate}
	
	\subsection*{For formal traceability questions}
	\begin{enumerate}[leftmargin=2em]
		\item \emph{Technical Appendix}
		\item \emph{Dependency Tables and Node Registry}
		\item \emph{Collected Propositions and Theorems}
		\item selected primary papers
	\end{enumerate}
	
	\subsection*{For review and audit use}
	\begin{enumerate}[leftmargin=2em]
		\item \emph{How to Verify the Theory}
		\item \emph{External Verification Checklist}
		\item \emph{Minimal Claims Register}
	\end{enumerate}
	
	\subsection*{For proof-oriented reading}
	\begin{enumerate}[leftmargin=2em]
		\item \emph{Collected Results and Proof Status Note}
		\item \emph{Collected Propositions and Theorems}
		\item \emph{Proof Status Register}
		\item \emph{Technical Appendix}
	\end{enumerate}
	
	\section{Internal structure of the package}
	
	The package may be understood as consisting of five coordinated layers.
	
	\subsection*{Layer 1: Terminology and interpretation}
	\begin{itemize}[leftmargin=2em]
		\item \emph{Glossary of Core Terms}
	\end{itemize}
	
	This layer controls how the main concepts should be read.
	
	\subsection*{Layer 2: Dependency logic}
	\begin{itemize}[leftmargin=2em]
		\item \emph{Acyclicity Statement and Dependency Criteria}
		\item \emph{Dependency Tables and Node Registry}
	\end{itemize}
	
	This layer controls how the series is structurally unfolded and how dependencies are assigned.
	
	\subsection*{Layer 3: Formal traceability}
	\begin{itemize}[leftmargin=2em]
		\item \emph{Technical Appendix}
	\end{itemize}
	
	This layer controls notation, assumptions, and formal prerequisite structure.
	
	\subsection*{Layer 4: Procedural and review support}
	\begin{itemize}[leftmargin=2em]
		\item \emph{How to Verify the Theory}
		\item \emph{External Verification Checklist}
	\end{itemize}
	
	This layer supports actual external use of the package.
	
	\subsection*{Layer 5: Claims, results, and proof status}
	\begin{itemize}[leftmargin=2em]
		\item \emph{Minimal Claims Register}
		\item \emph{Collected Results and Proof Status Note}
		\item \emph{Collected Propositions and Theorems}
		\item \emph{Proof Status Register}
	\end{itemize}
	
	This layer stabilizes what the theory centrally claims, what it has formally established, and
	how ready those results are for later proof-compendium treatment.
	
	\section{What this package does not try to do}
	
	The verification package is designed to support external verification, not to replace the series itself.
	
	It does \emph{not} attempt to:
	\begin{itemize}[leftmargin=2em]
		\item function as a substitute for the primary papers,
		\item eliminate the need for formal peer review,
		\item provide a full proof compendium,
		\item settle every possible interpretive dispute in advance,
		\item convert all architectural claims into theorem-grade results prematurely.
	\end{itemize}
	
	Its role is more practical: to make the existing theory easier to inspect, reconstruct, and
	verify from outside the original development process.
	
	\section{Current status of the package}
	
	At the present stage, the package should be understood as a complete verification-support layer
	for the current corpus of the \emph{General Theory of Cognitive Structuring} series.
	
	In particular, the package already provides:
	\begin{itemize}[leftmargin=2em]
		\item interpretive stabilization of core terms,
		\item explicit dependency criteria,
		\item a concrete node and paper registry,
		\item formal traceability for notation and assumptions,
		\item a procedural verification roadmap,
		\item a practical review checklist,
		\item a compact claims layer,
		\item a collected result layer,
		\item a collected proposition/theorem layer,
		\item and a proof-status layer.
	\end{itemize}
	
	This is sufficient to support informed external review of the theory's structure, terminology,
	dependency logic, and current level of formal consolidation.
	
	\section{Likely future extension}
	
	The most natural future extension of the package is a canonical \emph{Proof Notes} layer or a
	full \emph{Proof Compendium}, built on top of:
	\begin{itemize}[leftmargin=2em]
		\item the collected proposition/theorem layer,
		\item the proof-status register,
		\item and the technical appendix.
	\end{itemize}
	
	At the current stage, however, the verification package already serves as a strong public
	support structure for external verification of the series.
	
	\section{Closing note}
	
	The \emph{General Theory of Cognitive Structuring} is accompanied by multiple verification
	documents because different verification tasks require different forms of support. Terminology, dependency, traceability, review practice, claims stabilization, and proof status are related but not identical tasks.
	
	This overview is intended to make that division of labor explicit and to help external readers	enter the package in a controlled and efficient way.
	
\end{document}